\clearpage
\section{Datenaufbereitung TPS}
\label{sec:tps}
Dieses Kapitel beschreibt, wie die im Feldkurs erhobenen tachymetrischen Messungen
für die Netzausgleichung aufbereitet wurden.

\subsection{Messkonfiguration und Instrumente}
\label{sec:messkionfiguration und instrumente}
Die tachymetrischen Messungen wurdem mit drei Leica Nova MS60 (interne Nr. 5, 6 und 8)
und einem Leica Nova TS60 (interne Nr. 2) durchgeführt 

\subsection{Datenübertragung und -sicherung}
\label{sec:datenuebertragung und -sicherung}
Die Messdaten der gesamten Feldkampagne wurden redundant auf MS Teams und auf einer
externen Festplatte (G5 Nr. 4) gesichert. Die Rohdaten wurden unverändert und getrennt
von den bearbeiteten Daten abgelegt.

\subsection{Vorgehen in Leica Infinity}
\label{sec:vorgehen in leica infinity}
Die Aufbereitung erfolgte in Leica Infinity. Pro Station wurden folgende Schritte durchgeführt:
\begin{enumerate}
    \item Import der Messdaten
    \item Kontrolle der gemessenen Visuren auf Vollständigkeit und Plausibilität
    \item Prüfung der Satzmittel und Standardabweichungen
    \item Deaktivierung einzelner Sätze bei Auffälligkeiten (Kapitelxxx)
    \item Kontrolle und Anpassung der Meteokorrekturen
    \item Nachbearbeitung einzelner Stationen
\end{enumerate}

\subsection{Metekorrekturen}
\label{sec:metekorrekturen}
Bei der Stationierung wurden die Meteowerte am Standpunkt eingegeben. Das Netz
überspannt jedoch einen Höhenunterschied von rund 570 m, sodass diese Werte die
Verhältnisse entlang langer Visuren nur bedingt abbilden. Die Meteokorrekturen wurden
daher bei langen Visuren überprüft und angepasst. In Leica Infinity können dafür auch
die Meteowerte der Zielpunkte berücksichtigt werden. Beim Standpunkt 6003 wurden die
Meteokorrekturen zusätzlich angebracht

\subsection{Stationsspezifische Bearbeitung}
\label{sec:stationsspezifische_bearbeitung}
Die folgenden Abschnitte listen alle Stationen auf, bei denen Sätze deaktiviert oder
Daten angepasst wurden.

\subsubsection{Gebiet Brichen }
\label{sec:stationsspezifische_bearbeitung}




